\section{Introduction}
\label{sec:introduction}

\subsection{Problem Statement}

Data access remains one of the areas where the mismatch between the programming language and the storage model produces the greatest number of structural errors. In traditional database access layers, queries are described as strings and validated only at execution time. This means that incorrect column names, incompatible types, invalid \code{JOIN} constructs, or wrongly supplied parameters are discovered late --- often against a real database and only after additional manual testing. At the same time, switching from one backend to another usually requires rewriting the entire data-access layer.

Object-relational mapping is the classical way of connecting an application's object model with persistent storage~\cite{fowler2002}. Most ORM implementations, however, remain strongly tied to runtime mechanisms: SQL strings, code generation, annotations, macros, or dynamic metadata descriptions. In the context of C++, this introduces an additional tension. The language provides a powerful type system, templates, \code{constexpr} evaluation, and \code{concepts}, yet database libraries rarely exploit these mechanisms far enough for the compiler to become a first-class validator of query structure and backend compatibility.

A distinctive feature of the present system is that, despite the name \ORMName{} library, its architecture is not restricted to the relational model alone. Instead, it uses the same C++ data model and the same query IR as a logical layer over relational and non-relational systems. This raises a second research question: which parts of the ORM abstraction are genuinely universal, and which parts must remain explicitly constrained through a capability-based backend contract.

The third layer of the problem concerns execution itself. Even when the logical model and query IR are type-correct, the practical value of a modern data-access library also depends on whether it can integrate asynchronous execution without breaking its type guarantees and without forcing the user back into a callback-based style. The thesis is therefore organised around two complementary axes: compile-time modelling and asynchronous execution architecture.

\subsection{Objectives and Tasks}

The primary objective of the thesis is to present the design and implementation of a \ORMName{} library for C++ that:

\begin{enumerate}[label=\arabic*.]
    \item models data and queries through statically typed C++ constructs;
    \item allows \code{SELECT}, \code{INSERT}, \code{UPDATE}, and \code{DELETE} queries to be built as \code{constexpr} objects;
    \item separates logical query structure from backend-specific materialisation;
    \item supports both relational and non-relational connectors through \code{connector\_trait<DB>};
    \item uses compile-time capability tags to reject unsupported operations;
    \item demonstrates how the C++ model determines both the logical schema and the physical data structure for different backend families;
    \item realises a coroutine-based asynchronous layer that extends the synchronous connector model in a formal and type-safe way;
    \item formalises a contract for adding new connectors, including cases where they support native asynchronous execution;
    \item validates the architecture through systematic unit and integration tests and through empirical benchmark evaluation.
\end{enumerate}

To achieve that objective, the work performs the following tasks: analysis of existing solutions; definition of a static entity model; construction of a typed query IR; design of the connector layer and capability mechanism; construction of a coroutine-based async layer; implementation of connectors for different classes of databases and different execution models; investigation of migrations, prepared queries, thread safety, and async connection management; formulation of extensibility criteria; and development of a benchmarking methodology for comparing synchronous and asynchronous execution.

\subsection{Scope and Limitations}

The thesis treats the library as a modelling, query-construction, and execution layer. Its scope therefore includes the entity model, query IR, placeholder mechanism, prepared queries, migration model, connector trait architecture, capability gating, the coroutine-based async layer, and the concrete backend implementations present in the repository. A substantial part of the evaluation is the scenario analysis of SQLite, PostgreSQL, MySQL, MongoDB, Redis, Neo4j, and Cassandra, together with benchmarking of synchronous and asynchronous execution modes.

Out of immediate scope remain production-scale optimisations such as distributed connection management, backend-complete migration strategies, live benchmarking across a full matrix of production environments, and full exploitation of future C++26 reflection. These are discussed as future work in Chapter~\ref{sec:future_work}.

\subsection{Thesis Structure}

The thesis is organised as follows.

\textbf{Chapter~\ref{sec:existing_solutions}} reviews existing C++ solutions for database access and ORM, and positions the present work with respect to them.

\textbf{Chapter~\ref{sec:theoretical_background}} presents the theoretical foundations of \ORMName{} and the storage models required for the later relational and non-relational analysis.

\textbf{Chapter~\ref{sec:coroutines_async_theory}} introduces the theoretical apparatus of coroutines and asynchronous execution in C++, which is necessary for understanding the second major contribution of the library.

\textbf{Chapter~\ref{sec:architecture}} describes the compile-time ORM architecture of the library: entity declarations, the query IR, \code{connector\_trait}, capability mechanisms, and the place of migrations within the overall design.

\textbf{Chapter~\ref{sec:async_execution_architecture}} presents the async execution architecture and shows how the coroutine-based layer connects to the synchronous connector model, the thread-pool strategy, and connection management.

\textbf{Chapter~\ref{sec:backend_scenarios}} analyses how the same logical abstraction and different execution strategies are materialised across relational, document, key-value, graph, and wide-column backends.

\textbf{Chapter~\ref{sec:verification_evaluation}} brings together the verification evidence and the empirical benchmark evaluation.

\textbf{Chapter~\ref{sec:verification_extensibility}} formalises the contract for extending the library with new connectors, including capability and async requirements.

\textbf{Chapter~\ref{sec:future_work}} outlines the main limitations and future development directions, while \textbf{Chapter~\ref{sec:conclusion}} summarises the results and contributions.
